\chapter{Introducción}
\section{Descripción del problema}
Paz Sur es una empresa comercial ubicada en Cusco, Perú. Según la información proporcionada por el equipo, sus trabajadores reciben capacitaciones internas y de marcas proveedoras. Algunas actividades se vinculan a temporadas comerciales, como campaña escolar, Navidad y Día de la Madre. El registro de asistencia se realiza parcialmente de forma manual y puede transcribirse después a Excel.

Este procedimiento requiere organizar la capacitación, identificar participantes y conservar un historial consultable. No se han confirmado cantidades de trabajadores, sedes, capacitaciones, costos, tiempos de transcripción ni tasas de error. Por tanto, el informe no atribuye cifras o problemas adicionales a la empresa.

\section{Objetivo general}
Desarrollar progresivamente un sistema de gestión de capacitaciones que permita organizar, consultar y verificar información, con implementación ejecutable, persistencia, controles de acceso, pruebas y documentación defendible por el equipo.

\section{Objetivos específicos}
\begin{enumerate}
  \item Inspeccionar el repositorio y separar hechos, supuestos y propuestas.
  \item Formular requisitos verificables y un backlog candidato trazable.
  \item Diseñar una arquitectura coherente con los recursos de cuatro estudiantes.
  \item Construir un primer flujo de extremo a extremo de acceso, usuarios, roles y capacitaciones.
  \item Ejecutar verificaciones disponibles y registrar resultados sin fabricar evidencia.
  \item Preparar el diseño y las prioridades de las iteraciones posteriores.
\end{enumerate}

\section{Alcance y límites del informe}
El primer incremento aborda HU-01, HU-02, HU-03, HU-04 y HU-16 como dependencia de autenticación. El diseño contempla las 28 historias; no se declara implementado todo ese alcance. Las reglas de asistencia, evaluaciones y certificados aún requieren validación. Un certificado puede existir en determinadas capacitaciones y no se presupone obligatorio.

Las versiones V1 a V4 expresan objetivos de evolución. No se crean etiquetas retroactivas para simular una secuencia histórica. El estado de terminado depende también de revisión humana e integración, además de pruebas locales.

\section{Organización del documento}
Los capítulos presentan contexto, proceso, requisitos, arquitectura, planificación, implementación, pruebas, evolución estructural y resultados. Los anexos conservan la bitácora y las fuentes Mermaid. Los fragmentos de código se importan de archivos reales del repositorio para reducir divergencias entre informe e implementación.
